% BANKS-SOURCE: pdf=69 print=59 kind=problem-section id=chapter05-problems
\begin{banksproblems}

% BANKS-SOURCE: pdf=69 print=59 kind=problem id=5.1
\BanksProblem{5.1}{*}
In $D=4$, write the general solution of the Dirac equation,
\begin{equation*}
  (\ii\slashed\partial+m)\psi=0,
\end{equation*}
in terms of Fourier coefficients $a_i(\mathbf p)u_i(p)$ for positive-frequency solutions and
$b_i^\dagger(\mathbf p)v_i(p)$ for negative frequencies. $(\slashed p-m)u_i=0$ and $(\slashed p+m)v_i=0$. There are
two linearly independent solutions of each equation (see Appendix C). Show that
$a_i^\dagger$ and $b_i^\dagger$ transform like the creation operators of massive spin-$\frac12$
particles, if the field $\psi$ transforms in the $[2,1]\oplus[1,2]$ representation of the Lorentz group. Show
that the commutation or anti-commutation relation
\begin{equation*}
  [\psi_a(x,t),\psi_b^\dagger(y,t)]_\pm=\delta_{ab}\delta^{D-1}(\mathbf{x}-\mathbf{y})
\end{equation*}
implies bosonic or fermionic commutation relations for the creation and annihilation
operators. Show that only the fermionic choice leads to a local field, which
anti-commutes with itself and its conjugate, at space-like separation.

% BANKS-SOURCE: pdf=69 print=59 kind=problem id=5.2
\BanksProblem{5.2}{}
Starting from the Weyl representation of the Dirac matrices, find a change of basis
that makes all the matrices imaginary (Majorana representation) and another one
in which $\gamma^0$ is $\sigma_3\otimes1$ and $(\gamma_\mu)^\dagger=\gamma^0\gamma_\mu\gamma^0$.

% BANKS-SOURCE: pdf=69 print=59 kind=problem id=5.3
\BanksProblem{5.3}{*}
The most general action of the discrete symmetries $T$, $C$, and $P$ on mass
eigenstates of spin-$\frac12$ particles is (the $\eta$s are complex phases)
\begin{equation*}
  \begin{aligned}
  T^\dagger a^\dagger(\mathbf p,s)T&=\eta_T a^\dagger(-\mathbf p,-s),\\
  P^\dagger a^\dagger(\mathbf p,s)P&=\eta_P a^\dagger(-\mathbf p,s),\\
  C^\dagger a^\dagger(\mathbf p,s)C&=\eta_C b^\dagger(\mathbf p,s).
  \end{aligned}
\end{equation*}
Recall that the $T$ operator is anti-unitary, while the others are unitary. If the
anti-particle operator $b$ is not the same as that of the particle (i.e. the particle
carries a charge), then there are in principle independent phases $\bar\eta_{T,P}$ for the anti-
particle. Find the connections between these phases implied by a homogeneous
transformation law for local fields, and write down the transformation laws for
Dirac fields in the Weyl representation.

% BANKS-SOURCE: pdf=69 print=59 kind=problem id=5.4
\BanksProblem{5.4}{*}
In $D=4$, the Dirac anti-commutation relations imply that independent elements of the algebra generated by the Dirac
matrices are the antisymmetrized products $\gamma^{\mu_1\ldots\mu_k}$, where $k$ runs from 0 to 4. Rewrite the $k=3,4$ cases in terms of the
matrices $\gamma_\mu$ and $\gamma_5$. Use the results of the previous problem to characterize the
$T$, $C$, and $P$ transformation laws of all Dirac bilinears $\bar\psi\gamma^{\mu_1\ldots\mu_k}\psi$.

% BANKS-SOURCE: pdf=70 print=60 kind=problem id=5.5
\BanksProblem{5.5}{*}
Given a set of Weyl fields $\nu_a^i$, the most general non-derivative, bilinear term one
can add to the Lagrangian has the form
\begin{equation*}
  M_{ij}\epsilon^{ab}\nu_a^i\nu_b^j+\mathrm{h.c.}
\end{equation*}
Anti-commutativity of the fields and anti-symmetry of the Levi-Civita symbol
imply that $M$ is a symmetric complex matrix. Suppose that there is a $\mathrm{U}(1)$
transformation acting on the $\nu_a^i$, which commutes with $M$. Then the bilinear
can be written as a sum over subsets of fields with charges $\pm q$ under $\mathrm{U}(1)$,
\begin{equation*}
  \sum_q\nu_a^I(q)\nu_b^J(-q)m_{IJ}(q).
\end{equation*}
The part of this sum with $q=0$ is called a Majorana
mass term, while the rest of the bilinear is called a Dirac mass term. $m_{IJ}(q)$
is a matrix acting only on the subspace of fields with charge $q$. Show that by
doing independent unitary transformations on the charge $q$ and $-q$ fields (for
$q\ne0$) we can transform each $m_{IJ}(q)$ into a diagonal matrix with real positive
entries. Thus, the fermion kinetic term for Dirac fermion masses actually has a $\mathrm{U}(1)^{\text{No. of charges}}$
symmetry and preserves $T$, $C$, and $P$. In the context of the strong interactions these $\mathrm{U}(1)$ charges
are called quark flavors. They are not exactly preserved by weak interactions, but this
is a small effect. Flavor is an example of an accidental approximate symmetry. What can you say about the
Majorana mass matrix with regard to putting it in a canonical form, and its
transformation under discrete Lorentz transformations?

% BANKS-SOURCE: pdf=70 print=60 kind=problem id=5.6
\BanksProblem{5.6}{}
Derive the Källen--Lehmann spectral representation for a Dirac field $\psi_\alpha(x)$ and
for a conserved vector current operator $J_\mu(x)$ satisfying $\partial_\mu J^\mu=0$.

% BANKS-SOURCE: pdf=70 print=60 kind=problem id=5.7
\BanksProblem{5.7}{}
Solve the Dirac Green function equation
\begin{equation*}
  (-\ii\slashed\partial-e\slashed A-m)S(x,y)=\delta^D(x-y)
\end{equation*}
in a constant background electromagnetic field $A_\mu(x)=F_{\mu\nu}x^\nu$.

% BANKS-SOURCE: pdf=70 print=60 kind=problem id=5.8
\BanksProblem{5.8}{*}
In $D=4$, show that $\gamma_5=\ii\gamma^0\gamma^1\gamma^2\gamma^3$ and that
$[\gamma^\mu,\gamma^\nu]_+=-2\eta^{\mu\nu}$. Prove that $\gamma_5^2=1$ and that
$[\gamma_5,\gamma^\mu]_+=0$.

% BANKS-SOURCE: pdf=70 print=60 kind=problem id=5.9
\BanksProblem{5.9}{*}
Use the rules of Grassmann integration to prove the analog of Wick’s theorem
for fermions.

% BANKS-SOURCE: pdf=70 print=60 kind=problem id=5.10
\BanksProblem{5.10}{*}
Using the transformation $X_\pm=(\bar\eta\pm\eta)/2$, in the Grassmann integral formula
for the determinant of $M$, for the case in which $M$ is anti-symmetric, relate the
determinant and Pfaffian formulae and prove that $\det A=(\operatorname{Pf}A)^2$.

% BANKS-SOURCE: pdf=70 print=60 kind=problem id=5.11
\BanksProblem{5.11}{*}
In $D=4$, derive the Feynman rules for the Yukawa interaction
\begin{equation*}
  \Delta\mathcal{L}=\phi(g_S\bar\psi\psi+g_P\bar\psi\gamma_5\psi)
\end{equation*}
between a spin-zero field and a Dirac fermion.

% BANKS-SOURCE: pdf=70 print=60 kind=problem id=5.12
\BanksProblem{5.12}{*}
In $D=4$, prove all of the trace and contraction identities in Appendix C, and generalize
them to at least one higher power of Dirac matrices. Evaluate all of the traces
when an additional $\gamma_5$ is inserted, along with the Dirac matrices $\gamma^\mu$.

% BANKS-SOURCE: pdf=71 print=61 kind=problem id=5.13
\BanksProblem{5.13}{*}
Consider an interaction
\begin{equation*}
  \bar\psi\Gamma^A\psi\,g_A\phi^A+\frac12(\phi^A)^2,
\end{equation*}
where $\Gamma^A$ runs over the 16 independent anti-symmetrized Dirac products in $D=4$, $\phi^A$ is a tensor field of the appro-
priate kind to make a Lorentz-invariant product, and $g_A$ is the same for each
component of an irreducible Lorentz representation. The fields $\phi^A$ have no
kinetic terms. Solve for them algebraically, and show that you get the most
general four-fermion interaction for the $\psi$ field, which preserves the symmetry
$\psi\to\ee^{\ii\alpha}\psi$. This trick can be generalized to higher powers of $\psi$ and to inter-
actions that don’t preserve the symmetry. Using this trick, one can formally do the
fermion functional integral in terms of the Green function for the Dirac
operator $-\ii\slashed\partial-\phi_A\Gamma^A$ and the determinant of this operator. If $\psi$ were a
boson, we would get the inverse determinant instead of the determinant. Show
that this replacement implies the extra Feynman rule of $(-1)$ for each closed
fermion loop.

% BANKS-SOURCE: pdf=71 print=61 kind=problem id=5.14
\BanksProblem{5.14}{*}
With the mostly-plus convention and positive-frequency modes proportional to
$\ee^{+\ii p\cdot x}$, prove the Gordon identity
\begin{equation*}
  2m\bar u(k)\gamma^\mu u(p)=\bar u(k)\left[-(p+k)^\mu+\gamma^{\mu\nu}(k-p)_\nu\right]u(p)
\end{equation*}
and the analogous identity for the anti-particle spinors $v(p)$.

% BANKS-SOURCE: pdf=71 print=61 kind=problem id=5.15
\BanksProblem{5.15}{*}
In $D=4$, solve the momentum-space Dirac equations
$(\slashed p-m)u(p,s)=0=(\slashed p+m)v(p,s)$,
in both the Dirac and the Weyl bases for the Dirac matrices (see Appendix C).

% BANKS-SOURCE: pdf=71 print=61 kind=problem id=5.16
\BanksProblem{5.16}{*}
Show that charge-conjugation symmetry implies that the representation of the
internal symmetry group $G$ is real or pseudo-real.

\end{banksproblems}
